\documentclass{article}
\input{../packages.tex}
\begin{document}
    \begin{titlepage}
        \begin{center}

        Міністерство освіти і науки України
        
        НТУУ «Київський політехнічний інститут»
        
        Фізико-технічний інститут
        \vspace{3.3cm}
        
        {\textbf{Блокчейн та основи криптовалют}\\Конспект 1. Вступна лекція.}

        \vspace{8.5cm}

        \begin{flushright}
            \textbf{Виконав:}\\Студент 4-го курсу\\групи ФІ-21\\Климентьєв Максим\\
            \textbf{Перевірив:}\\\text{\_\_\_\_\_\_\_\_\_\_\_\_\_\_\_\_\_\_}
        \end{flushright}

        \end{center}
    \end{titlepage}
    \newpage

    \pagenumbering{gobble}
    \tableofcontents
    \cleardoublepage
    \pagenumbering{arabic}
    \setcounter{page}{3}

    \newpage
    \section{Криптографічні Геш-функції}
    Геш-функція - функція, яка приймає бітову послідовність скінченної довжини, а видає на виході фіксованої довжини (256 біт зазвичай).

    Криптографічна геш-функція - це геш-функція, в якій:
    \begin{itemize}
        \item відсутні колізії (Якщо обчислювально складно знайти такі $ x \neq y: H(x) = H(y) $, якщо перебрати приблизно $2^130$ входів то з ймовірністю $99.8\%$ можна знайти колізію),
        \item сховано прообраз (Маючи $H(x)$ ми не можемо знайти використаний $x$, зазвичай складніше - навіть якщо є $r$, який обирається рівноймовірно з великої множини, і якщо відомо для багатьох $r$ $H(r \vert x)$), ми не зможемо знайти використаний $x$ ,
        \item функція є дружньою з головоломками.
    \end{itemize}

    Геш-функція корисна для:
    \begin{itemize}
        \item розпізнавання (порівняння) файлів,
        \item фіксування значення за допомогою геш-функції так, що його неможливо дізнатись раніше часу та неможливо змінити після розкриття, оскільки спочатку публікується геш, а згодом саме значення та сіль, і через стійкість геш-функції до колізій підміна значення є обчислювально неможливою,
    \end{itemize}

\end{document}
